\section{Calibración de las tareas para el modelo objetivo}
\label{sec:calibracion}

\guia{Usar la línea base para evaluar si las tareas ofrecen una zona
aprendible para el modelo objetivo. El enunciado pide evitar tanto
el fracaso sistemático como tareas que el modelo siempre resuelve;
no establece umbrales numéricos obligatorios.}

\subsection{Evidencia por nivel de dificultad}
\pendiente{interpretar conjuntamente \passuno\ y \passk\ del modelo
objetivo en cada nivel. Identificar familias de tareas con aciertos,
margen de mejora o ausencia de aciertos en la muestra evaluada.}

\subsection{Ajustes propuestos y alcance de la evidencia}
\pendiente{justificar qué tareas conservar, simplificar, endurecer o
diversificar, y cómo se verificaría el efecto de esos cambios.}

\guia{Precisar si los ajustes se propusieron o se evaluaron realmente.
Utilizar validación para calibrar y reservar prueba para la evaluación
final. Los resultados de inferencia orientan la calibración; por sí
solos no demuestran una mejora futura mediante entrenamiento.}
